\subsection{Trajectory Smoothing and Stop Avoidance}
% EN: Review velocity planning, energy-aware control, control barrier functions, and smooth crossing strategies.
% CN: 回顾速度规划、能耗优化、控制屏障函数和平滑通过路口的方法。
\begin{addedtodayblock}
Recent work shows that intersection efficiency depends not only on the discrete passing order but also on whether the assigned order can be executed by dynamically feasible vehicle trajectories. Graph-based CAV cooperation methods model crossing, diverging, converging, and reachability conflicts before selecting a passing order, which is close in spirit to representing conflict zones as shared scheduling resources \cite{chen2021conflictfree,chen2021lanechanging}. Priority-aware resequencing further shows that strict FCFS/FIFO decision orders can be relaxed through replanning while still passing target times to a lower-level optimal controller \cite{chalaki2021priority}. These studies motivate the scheduling layer in this paper, but they usually rely on handcrafted graph algorithms or analytic resequencing rather than an offline-trained graph policy.

Trajectory-level studies address the complementary question of how a vehicle should move after a desired crossing time, phase, or maneuver has been selected. Eco-driving and eco-approach controllers use reinforcement learning, Pontryagin-minimum-principle solutions, or hybrid rule-learning structures to reduce energy consumption, delay, and unnecessary acceleration near intersections \cite{bai2022hybrid,jiang2023ecodriving,lakshmanan2022cooperative,dong2023overtaking}. These papers support the use of stop count, acceleration variation, and energy proxy metrics in the evaluation, because stop-line waiting is not merely a delay phenomenon but also changes the smoothness and energy profile of the vehicle trajectory.

Safety-oriented trajectory planning adds constraints that the high-level schedule alone cannot guarantee. Predictive-control and spatial-domain formulations derive trajectories under rear-end, crossing, merging, diverging, and mixed-traffic uncertainty constraints \cite{mahbub2022safety,zhao2025spatial}. Learning-based AIM approaches can improve flexibility and travel time, but recent work also emphasizes that learned decisions need explicit safety evaluation or filtering, such as uncertainty-aware high-order control barrier functions \cite{luo2023realtime,yu2025uncertainty,cederle2024distributed,li2023dhal}. A broad CAV-control survey similarly identifies trajectory tracking, collaborative control, safety, comfort, and energy saving as coupled requirements rather than independent design goals \cite{liu2023systematic}.

The proposed framework follows this schedule-to-trajectory view. The high-level JSSP/RL scheduler determines the conflict-zone order and target entry times, while the low-level smoother treats those times as references to be tracked under speed, acceleration, jerk, and safety constraints. If the smoother cannot track the assigned reservations without violating collision-avoidance or comfort bounds, it reports infeasibility and triggers replanning or conservative stop-line fallback. This feedback interface is the main distinction from methods that optimize only the passing order or only the vehicle trajectory in isolation.
\end{addedtodayblock}
